\documentclass[conference]{IEEEtran}
\IEEEoverridecommandlockouts
\usepackage{amsmath,amssymb,amsfonts}
\usepackage{graphicx}
\usepackage{textcomp}
\usepackage{xcolor}
\usepackage{url}
\usepackage{hyperref}
\def\BibTeX{{\rm B\kern-.05em{\sc i\kern-.025em b}\kern-.08em
    T\kern-.1667em\lower.7ex\hbox{E}\kern-.125emX}}
\begin{document}

\title{Tarea 03: Generación de imágenes con autoencoder\\
PyTorch Lightning -- Hydra -- WandB}

\author{\IEEEauthorblockN{Fabricio Mena Mejía}
\IEEEauthorblockA{Instituto Tecnológico de Costa Rica\\
Email: fabriciomena11@estudiantec.cr}
\and
\IEEEauthorblockN{Anthony Fuentes Calvo}
\IEEEauthorblockA{Instituto Tecnológico de Costa Rica\\
Email: anthonyfuentescalvo@estudiantec.cr}
\and
\IEEEauthorblockN{Fernando Sánchez Hidalgo}
\IEEEauthorblockA{Instituto Tecnológico de Costa Rica\\
Email: Fer.sanchez@estudiantec.cr}}

\maketitle

\begin{abstract}
Este documento describe el desarrollo de la Tarea~03 del curso de Inteligencia Artificial: la implementación y evaluación de modelos generativos basados en autoencoders para la reconstrucción de imágenes industriales. Siguiendo el enunciado, se experimentan dos arquitecturas---un Variational Autoencoder (VAE) y un autoencoder tipo U-Net con conexiones salto---sobre el dataset MVTec AD, restringido a las clases \textit{cable}, \textit{capsule}, \textit{screw} y \textit{transistor}, con imágenes de entrada de $128 \times 128$ píxeles y tres canales. El entrenamiento se estructura con PyTorch Lightning, la gestión modular de hiperparámetros con Hydra y el registro comparativo de experimentos en Weights \& Biases. Para cada arquitectura se prevén al menos cuatro corridas que difieren en la función de pérdida (L1, L2, SSIM y SSIM+L1), con el fin de analizar reconstrucciones en validación y prueba, el comportamiento del espacio latente (t-SNE) y la separabilidad del error de reconstrucción entre muestras normales y anómalas.
\end{abstract}

\begin{IEEEkeywords}
Autoencoder, variational autoencoder, U-Net, detección de anomalías, MVTec AD, PyTorch Lightning, Hydra, Weights \& Biases, reconstrucción de imágenes.
\end{IEEEkeywords}

\section{Introducción}

La detección de defectos en entornos industriales puede abordarse de forma no supervisada aprendiendo a reconstruir únicamente patrones normales de producto o textura. Cuando un autoencoder se entrena con imágenes sin anomalías, las reconstrucciones tienden a ser fieles en casos normales y a degradarse ante defectos no vistos, lo que convierte el error de reconstrucción en una señal útil para la detección.

La Tarea~03 del curso de Inteligencia Artificial solicita implementar y comparar dos familias de modelos---VAE y U-Net---sobre el benchmark MVTec AD, utilizando un subconjunto de cuatro clases y resolución fija de $128 \times 128$. El flujo experimental debe ser reproducible mediante configuraciones Hydra, el ciclo de entrenamiento debe organizarse con PyTorch Lightning, y los resultados (pérdidas, reconstrucciones y visualizaciones del espacio latente) deben centralizarse en WandB para su comparación sistemática entre funciones de pérdida.

Este informe acompaña el notebook de la tarea y documentará, en secciones posteriores, el dataset y preprocesamiento, el diseño de arquitecturas, los experimentos de entrenamiento y el análisis crítico de reconstrucciones en imágenes normales y defectuosas.

\end{document}
